\chapter{Exercises}
This chapter holds a selection of a few general exercises that we believe hold high didactic caliber, separated by chapter. Many of these exercises aim to directly show the utility of the theory developed until this point, and our personal understanding on the subject was aided through the solution of the proposed exercises. 

They do not, however, aim to prepare for any sort of test on the subject, and the texts of the following exercises are to be taken as didactic tools rather than exam questions.


\clearpage
\section{Probability 101}\label{exer:probability_101}

The following exercises will be focused on performing variable change and finding the distribution of various statistics.

\begin{exercise}\label{exer:variable_sum}
	Given two random variables $X$ and $Y$ with probability density functions
	$p_X(x)$ and $p_Y(y)$, respectively, we want to determine the pdf
	$f_S(s)$ of the statistic
	\[
	S = X + Y
	\]
	We consider two separate cases:
	\begin{itemize}
		\item $X$ and $Y$ dependent
		\item $X$ and $Y$ independent
	\end{itemize}
	
	The main strategy is to express one of the original variables as a function
	of the statistic, perform a change of variables, and then marginalize over
	the remaining variable:
	\[
	(x,y)\longrightarrow (x,s=x+y)
	\]
	
	For a general transformation
	\[
	(x_1,\dots,x_N)\longrightarrow(y_1,\dots,y_N)
	\]
	the probability density transforms as
	\[
	p_Y(y_1,\dots,y_N)
	=
	p_X(x_1,\dots,x_N)
	\left|
	\det\!\left(
	\frac{\partial y_i}{\partial x_j}
	\right)
	\right|^{-1}_{x=x(y)}
	\]
	
	In our case, $y=s-x$, the Jacobian matrix is
	\[
	\begin{pmatrix}
		\dfrac{\partial x}{\partial x} &
		\dfrac{\partial x}{\partial y}
		\\[1em]
		\dfrac{\partial s}{\partial x} &
		\dfrac{\partial s}{\partial y}
	\end{pmatrix}
	=
	\begin{pmatrix}
		1 & 0\\
		1 & 1
	\end{pmatrix}
	\]
	whose determinant is
	\[
	J = 1
	\]
	
	Thus, the joint pdf of $(X,S)$ is
	\[
	p_{X,S}(x,s)
	=
	p_{X,Y}(x,s-x)\frac{1}{|J|}
	=
	p_{X,Y}(x,s-x)
	\]
	
	To obtain the distribution of $S$, we marginalize over $x$:
	\[
	f_S(s)
	=
	\int_{-\infty}^{+\infty}
	p_{X,Y}(x,s-x)\,dx
	\]
	
	If $X$ and $Y$ are independent, then
	\[
	p_{X,Y}(x,y)=p_X(x)p_Y(y)
	\]
	and therefore
	\[
	f_S(s)
	=
	\int_{-\infty}^{+\infty}
	p_X(x)p_Y(s-x)\,dx
	\]
	
	This integral is called the \emph{convolution} of the two \emph{pdfs}:
	\[
	f_S = p_X * p_Y
	\]
	
	Hence, the sum of two independent random variables is distributed
	according to the convolution of their probability densities.
	
	It is important to note that one cannot simply add the two \emph{pdfs},
	i.e.\ write
	\[
	f_S(s)\neq p_X(s)+p_Y(s)
	\]
	This is because a pdf describes how probability mass is distributed
	over the possible values of a random variable. The event
	$S=s$ can be obtained through infinitely many combinations
	\[
	x+(s-x)=s
	\]
	so all possible decompositions of $s$ must be taken into account.
	The convolution precisely performs this weighted sum over all
	possible pairs $(x,s-x)$.
	
	Consider now the case where both $X$ and $Y$ follows a uniform distribution between $0$ and $m$ $\mathcal{U}[0,m]$. We have that: $s \in [0, 2m]$.
	
	We have then:
	\[
	f_S(s)
	=
	\int_{0}^{0}
	\frac{1}{m}p_Y(s-x)\,dx
	\]
	knowing that $p_Y(s-x)\neq0 \; \text{if} s-x\leq m$ we must separate the domain of $f_S(s)$. Thus:
	\[
		f_S(s)=\begin{cases}
		\int_0^s\frac{1}{m^2}\;dx=\frac{s}{m^2} \; &\text{if} \; 0\leq s\leq m\\
		\int_{s-m}^m\frac{1}{m^2}\;dx=\frac{2m-s}{m^2} \; &\text{if} \; m\leq s\leq 2m\\
		0 \; &\text{otherwise}
		\end{cases}
	\]
	This shows that the density increases linearly up to $s=m$ and then decreases symmetrically. This is expected: there are more ways to obtain intermediate sums than extreme ones.
	
	We can verify normalization:
	\[
		\int_{0}^{2m}f_S(s)\;ds=\int_0^m\frac{s}{m^2}\;dx+\int_m^{2m}\frac{2m-s}{m^2}\;ds=\frac{5}{2}-\frac{3}{2}=1
	\]
	
	We are now interested in computing the expected value and variance. For the expected value:
	\[
		\E[s]=\E[x+y]=\E[x]+\E[y]=\frac{m}{2}+\frac{m}{2}=\frac{m}{2}
	\]
	where we've used the nice properties of the expected value. For the variance:
	\begin{align*}
		\Var[s]&=\E[(x+y-m)^2]=\E\left[\left(x-\frac{m}{2}+y-\frac{m}{2}\right)^2\right]\\
		&=\underbrace{\E\left[\left(x-\frac{m}{2}\right)^2\right]}_{\Var[x]}+\underbrace{E\left[\left(y-\frac{m}{2}\right)\right]}_{\Var[y]}+2\underbrace{E\left[\left((x-\frac{m}{2})\right)\left(y-\frac{m}{2}\right)\right]}_{\Cov[x,y]}\\
		&=\frac{m^2}{16}+\frac{m^2}{16}+0=\frac{m^2}{6}
	\end{align*}
	where the covariance is null since $x$ and $y$ independent. The reader can also compute expected value and variance directly from the pdf to become comfortable with these types of calculations.
\end{exercise}

\clearpage
\begin{exercise}
	Following the same blueprint as exercise \ref{exer:variable_sum}, we now want to determine the pdf
	$f_{\Pi}(\pi)$ of the statistic
	\[
	\Pi = X \cdot Y
	\]
	As we've done before, we want to express one of the original variables as a function
	of the statistic, perform a change of variables, and then marginalize over
	the remaining variable:
	\[
	(x,y)\longrightarrow (x,\pi=xy)
	\]
	For the Jacobian:
	\[
		J = \det 	\begin{pmatrix}
			\dfrac{\partial x}{\partial x} &
			\dfrac{\partial x}{\partial y}
			\\[1em]
			\dfrac{\partial \pi}{\partial x} &
			\dfrac{\partial \pi}{\partial y}
		\end{pmatrix} = \det \begin{pmatrix}
		1 & 0\\
		y & x \\
		\end{pmatrix}=x
	\]
	Thus:
	\[
		p_{X, \Pi}(x, \pi)=p_{X, Y}\left(x, \frac{\pi}{x}\right)\frac{1}{\left|x\right|}
	\]
	Marginalizing over $x$:
	\[
		f_{\Pi}(\pi)=\int_{-\infty}^{+\infty}p_{X, Y}\left(x, \frac{\pi}{x}\right)\frac{1}{\left|x\right|}\; dx
	\]
	that for independent variables:
	\[
		f_{\Pi}(\pi)=\int_{-\infty}^{+\infty}p_{X}(x)p_{Y}\left(\frac{\pi}{x}\right)\frac{1}{\left|x\right|}\; dx
	\]
\end{exercise}

\clearpage
\begin{exercise}
	Now you know the drill. We want to determine the pdf $f_{R}(r)$ of the statistic:
	\[
		R=\frac{X}{Y}
	\]
	We set \sidenotemark:
	\[
		(x,y)\longrightarrow (r=x/y,y)
	\]
	Thus:
	\[
	J = \det 	\begin{pmatrix}
		\dfrac{\partial r}{\partial x} &
		\dfrac{\partial r}{\partial y}
		\\[1em]
		\dfrac{\partial x}{\partial x} &
		\dfrac{\partial x}{\partial y}
	\end{pmatrix} = \det \begin{pmatrix}
		\frac{1}{y} & -\frac{x}{y^2}\\
		0 & 1\\
	\end{pmatrix}=\frac{1}{y}
	\]
	We have that:
	\[
		p_{R, Y}(r,y)=p_{X,Y}(ry,y)|y|
	\]
	marginalizing over $y$:
	\[
		f_R(r)=\int_{-\infty}^{+\infty}p_{X,Y}(ry,y)|y|\;dy
	\]
	that, for independent variables, becomes:
	\[
		f_R(r)=\int_{-\infty}^{+\infty}p_X(rx)p_Y(y)|y|\;dy
	\]
\end{exercise}
\sidenotetext{This time we substituted the variable $x$ with the statistic $r$. Obviously there is no difference in the final result, but the reader can easily verify that this choice greatly simplifies the algebra.}

\clearpage
\section{Functional Inferences}\label{exer:functional_inf}

\clearpage
\section{Point Estimations}\label{exer:point_est}

During the critical years of 1941 and 1942, Allied command faced a daunting question: exactly how many formidable Mark IV and V Panzer tanks was Germany producing? Conventional military intelligence relied on traditional espionage: spies, reconnaissance, and battlefield counts. However, these methods yielded wildly contradictory and inflated estimates, suggesting a staggering production rate of around $1,400$ tanks per month \cite{The_Guardian_2006}. Facing the daunting prospect of a continental invasion, the Allies needed certainty, so they turned to an unconventional weapon: statistical intelligence. 

The breakthrough came from analyzing the serial numbers on captured German tanks. They noticed that the Germans, with characteristic bureaucratic precision, stamped a serial number on almost every component, and assumed that those serial numbers were sequential. Using a remarkably small sample, the Allies were able to estimate a production rate of $246$ tanks per month between June 1940 and September 1942. When post-war records were finally uncovered, the actual German production logs revealed the true number to have been $245$ per month. The statisticians had achieved near-perfect accuracy, proving that a handful of sequential serial numbers could out-spy an entire intelligence network.

In the exercise below, we will step into the shoes of those Allied statisticians and give our own estimate for the maximum of a sample of uniformly-distributed numbers.

\begin{exercise}
	We assume the serial numbers of the captured tanks to be uniformly distributed integers between $1$ and $N$, where $N$ is the total number of tanks produced, which we aim to estimate. Let's use the \emph{MLE}, which we know to be (see appendix \ref{sec:uniform}), for a uniform distribution, $\hat{N}_{MLE} = \max{x_i}$, where $\{x_i\}$ are the individual serial numbers of the captured tanks. Let's assume to have captured and analyzed $k$ tanks.
	
	It's important to know the bias of this estimate, as the fate of important strategic decisions is on the line and we would not want to over or under estimate the number of tanks produces. Once again dipping into appendix \ref{sec:uniform} we see that the expected value for $x_{max}$ is:
	
	\[
	\E\left[x_{max}\right] = N \frac{k}{k+1} \Rightarrow b\left[{N}_{MLE}\right] = -N\frac{1}{k+1}
	\]
	
	Since the Allies only captured a few tanks, an asymptotically unbiased estimate was not good enough, and we build a new, unbiased estimate of $N$:
	
	\[
	\hat{N}^\prime = \frac{k+1}{k} x_{max}
	\]
	
	Which proves to be an efficient estimator for the number of tanks produced. This exercise, though simple, illustrates an important historical usage of the concepts learned in this course.	
\end{exercise}


\clearpage
PROBLEMA 3 COMPITO 10/01/18 PUNTO 3

\clearpage
\section{Interval Estimations}\label{exer:interval_est}

\clearpage
\section{Hypothesis Testing}\label{exer:hypoythesis_test}

PROBLEMA 3 COMPITO 10/01/18 PUNTI 1 E 2

\clearpage
\section{Goodness of Fit}\label{exer:gof}